\documentclass[11pt]{article}
\usepackage[a4paper,margin=1in]{geometry}
\usepackage{fontspec}
\usepackage[boldfont=false]{xeCJK}
\setCJKmainfont{FandolSong-Regular.otf}
\usepackage{amsmath}
\usepackage{amssymb}
\usepackage{array}
\usepackage{booktabs}
\usepackage{graphicx}
\usepackage{adjustbox}
\usepackage{enumitem}
\setlist[itemize]{topsep=2pt,partopsep=0pt,parsep=0pt,itemsep=2pt,leftmargin=1.4em}
\setlist[enumerate]{topsep=2pt,partopsep=0pt,parsep=0pt,itemsep=2pt,leftmargin=1.6em}
\usepackage{parskip}
\usepackage{fvextra}
\fvset{breaklines=true,breakanywhere=true,fontsize=\small}
\setlength{\parindent}{0pt}

\begin{document}

\begin{center}
  {\LARGE\bfseries Chemical energetics}\\[0.35em]
  {\large A Level Chemistry}
\end{center}
\vspace{0.6em}

\section*{Lattice energy and Born–Haber cycles}

These cycles use several \textbf{enthalpy changes} 焓变 ($\Delta H$). Two new ones are:

\begin{itemize}
\item the \textbf{enthalpy change of atomisation} 原子化焓变, $\Delta H_{\text{at}}$ — the energy to make \textbf{one mole} of gaseous atoms from an element. It is always positive (bonds must break).

\item the \textbf{lattice energy} 晶格能, $\Delta H_{\text{latt}}$ — the energy change when one mole of a solid ionic lattice forms from its gaseous ions. It is always negative (strong bonds form).

\end{itemize}

\subsection*{Electron affinity}

The \textbf{first electron affinity} 电子亲和能 (EA) is the energy change when one mole of gaseous atoms each gain one electron to form one mole of $1-$ ions. The first EA is usually negative.

The same factors as ionisation energy apply (nuclear charge, atomic radius, shielding). Going down Group 16 or 17, the EA becomes \textbf{less} exothermic, because the atom is larger and pulls the extra electron in less strongly. (The very top element is an exception: its atom is so small that electron repulsion makes its EA less exothermic than the one below it.)

\subsection*{Born–Haber cycles}

A \textbf{Born–Haber cycle} 玻恩哈伯循环 is an energy cycle that links the enthalpy change of formation of an ionic solid with its atomisation, ionisation energy, electron affinity and lattice energy. Using Hess's law, you go round the cycle to find any one unknown step. (You only need $+1$ and $+2$ cations and $-1$ and $-2$ anions.)

\subsection*{What controls the size of a lattice energy}

The lattice energy is more negative (stronger) when:

\begin{itemize}
\item the \textbf{ionic charge} is higher (e.g. $\text{Mg}^{2+}$ beats $\text{Na}^+$).

\item the \textbf{ionic radius} is smaller.

\end{itemize}

Both make the attraction between the ions stronger.

\section*{Enthalpies of solution and hydration}

\begin{itemize}
\item the \textbf{enthalpy change of hydration} 水合焓变, $\Delta H_{\text{hyd}}$ — the energy change when one mole of gaseous ions is surrounded by water to form aqueous ions. It is exothermic.

\item the \textbf{enthalpy change of solution} 溶解焓变, $\Delta H_{\text{sol}}$ — the energy change when one mole of solute dissolves fully in water.

\end{itemize}

An energy cycle links the three:

\[
\Delta H_{\text{sol}} = -\Delta H_{\text{latt}} + \Delta H_{\text{hyd}}
\]

(To dissolve, you first pull the lattice apart, then hydrate the ions.) Like lattice energy, $\Delta H_{\text{hyd}}$ is more exothermic for ions with a higher charge and a smaller radius.

\section*{Entropy change, ΔS}

\textbf{Entropy} 熵 ($S$) measures the number of ways the particles and their energy can be arranged in a system. More ways means more "disorder".

The \textbf{entropy change} 熵变 ($\Delta S$) is \textbf{positive} when disorder increases, and \textbf{negative} when it falls:

\begin{center}
\adjustbox{max width=\linewidth}{%
\begin{tabular}{ll}
\toprule
\textbf{Change} & \textbf{Sign of $\Delta S$} \\
\midrule
solid → liquid → gas (melting, boiling), or dissolving & positive \\
a rise in temperature & positive \\
a reaction that makes \textbf{more} gas molecules & positive \\
a reaction that makes \textbf{fewer} gas molecules & negative \\
\bottomrule
\end{tabular}}
\end{center}

You can calculate it from standard entropies:

\[
\Delta S^{\ominus} = \sum S^{\ominus}(\text{products}) - \sum S^{\ominus}(\text{reactants})
\]

\section*{Gibbs free energy change, ΔG}

The \textbf{Gibbs free energy} 吉布斯自由能 change decides whether a reaction can happen on its own:

\[
\Delta G^{\ominus} = \Delta H^{\ominus} - T\Delta S^{\ominus}
\]

Here $T$ is the temperature in kelvin. A reaction is \textbf{feasible} (it can happen) when $\Delta G$ is negative or zero. The \textbf{feasibility} 可行性 therefore depends on temperature:

\begin{center}
\adjustbox{max width=\linewidth}{%
\begin{tabular}{lll}
\toprule
\textbf{$\Delta H$} & \textbf{$\Delta S$} & \textbf{When feasible} \\
\midrule
negative & positive & at all temperatures \\
positive & positive & only at high temperature \\
negative & negative & only at low temperature \\
positive & negative & never \\
\bottomrule
\end{tabular}}
\end{center}

To find the changeover temperature, set $\Delta G = 0$, which gives $T = \Delta H / \Delta S$.


\end{document}
